\documentclass[11pt,a4paper]{article}

\usepackage[utf8]{inputenc}
\usepackage[T1]{fontenc}
\usepackage[spanish]{babel}
\usepackage{amsmath,amssymb,mathtools}
\usepackage{geometry}
\usepackage{lmodern}
\usepackage{microtype}

\geometry{margin=2.4cm}
\setlength{\parindent}{0pt}
\setlength{\parskip}{0.45em}

\newcommand{\R}{\mathbb{R}}
\newcommand{\sen}{\operatorname{sen}}
\newcommand{\rg}{\operatorname{rg}}
\newcommand{\tr}{\operatorname{tr}}

\begin{document}

\begin{center}
{\LARGE\bfseries Geometría Lineal}\\[0.4em]
{\Large Tema 5: Cónicas y cuádricas}\\[0.4em]
{\large Soluciones}
\end{center}

\section{Ejercicio 1}
\subsection{Solución}

La cónica es:

\[
x^2-4xy+3y^2-2x-y-5=0.
\]

Su matriz homogénea es:

\[
\boxed{
A=
\begin{pmatrix}
1&-2&-1\\
-2&3&-\frac12\\
-1&-\frac12&-5
\end{pmatrix}
}
\]

y:

\[
\boxed{
\widetilde X^tA\widetilde X=0,
\qquad
\widetilde X=
\begin{pmatrix}
x\\y\\1
\end{pmatrix}.
}
\]

Equivalentemente:

\[
A_0=
\begin{pmatrix}
1&-2\\
-2&3
\end{pmatrix},
\qquad
L=
\begin{pmatrix}
-1&-\frac12
\end{pmatrix},
\qquad
C_1=-5.
\]

Por tanto:

\[
\boxed{
X^tA_0X+2LX+C_1=0.
}
\]

\section{Ejercicio 2}
\subsection{Solución}

Tomamos la ecuación general:

\[
ax^2+2bxy+cy^2+2dx+2ey+f=0.
\]

Sustituyendo los cinco puntos se obtiene:

\[
a=\frac e4,
\qquad
b=-\frac e2,
\qquad
c=e,
\qquad
d=-\frac e4,
\qquad
f=0.
\]

Tomando:

\[
e=4,
\]

se obtiene:

\[
\boxed{
x^2-4xy+4y^2-2x+8y=0.
}
\]

\section{Ejercicio 3}
\subsection{Solución}

Los focos son:

\[
F_1=(1,-2),
\qquad
F_2=(1,4).
\]

El centro es:

\[
C=(1,1)
\]

y:

\[
c=3.
\]

Para el punto:

\[
P=(5,1),
\]

se tiene:

\[
PF_1=PF_2=5.
\]

Por tanto:

\[
2a=10,
\qquad
a=5.
\]

Además:

\[
b^2=a^2-c^2=25-9=16.
\]

Como el eje mayor es vertical:

\[
\boxed{
\frac{(x-1)^2}{16}
+
\frac{(y-1)^2}{25}
=
1.
}
\]

\section{Ejercicio 4}
\subsection{Solución}

El vértice es:

\[
V=(3,2)
\]

y el foco:

\[
F=(5,2).
\]

Por tanto:

\[
p=2.
\]

La parábola tiene eje horizontal y abre hacia la derecha.

Así:

\[
\boxed{
(y-2)^2=8(x-3).
}
\]

\section{Ejercicio 5}
\subsection{Solución}

La ecuación es:

\[
(x+6)^2=-24(y-2).
\]

Comparando con:

\[
(x-h)^2=4p(y-k),
\]

se obtiene:

\[
h=-6,
\qquad
k=2,
\qquad
p=-6.
\]

Por tanto:

\[
\boxed{
V=(-6,2)
}
\]

\[
\boxed{
F=(-6,-4)
}
\]

y la directriz es:

\[
\boxed{
y=8.
}
\]

\section{Ejercicio 6}
\subsection{Solución}

La hipérbola es:

\[
x^2-\frac{y^2}{9}=1.
\]

Así:

\[
a=1,
\qquad
b=3,
\qquad
c=\sqrt{a^2+b^2}=\sqrt{10}.
\]

Los vértices son:

\[
\boxed{
(\pm1,0)
}
\]

y los focos:

\[
\boxed{
(\pm\sqrt{10},0).
}
\]

Las asíntotas son:

\[
\boxed{
y=\pm3x.
}
\]

\section{Ejercicio 7}
\subsection{Solución}

La hipérbola es:

\[
\frac{(y-2)^2}{4}
-
\frac{(x+2)^2}{4}
=
1.
\]

Por tanto:

\[
a=2,
\qquad
b=2,
\qquad
c=2\sqrt2.
\]

Los vértices son:

\[
\boxed{
(-2,4),
\qquad
(-2,0)
}
\]

y los focos:

\[
\boxed{
(-2,2+2\sqrt2),
\qquad
(-2,2-2\sqrt2).
}
\]

Las asíntotas son:

\[
\boxed{
y-2=\pm(x+2).
}
\]

\section{Ejercicio 8}
\subsection{Solución}

\textbf{a)}

\[
x^2+2xy+y^2+2x-1=0.
\]

La matriz cuadrática es:

\[
A_0=
\begin{pmatrix}
1&1\\
1&1
\end{pmatrix}
\]

y:

\[
|A_0|=0,
\qquad
|A|\neq0.
\]

Por tanto:

\[
\boxed{\text{Parábola}.}
\]

\textbf{b)}

\[
2xy+2x-2y+1=0.
\]

Se tiene:

\[
A_0=
\begin{pmatrix}
0&1\\
1&0
\end{pmatrix},
\qquad
|A_0|=-1<0.
\]

Por tanto:

\[
\boxed{\text{Hipérbola}.}
\]

\textbf{c)}

\[
x^2+y^2-2x+2y-3=0.
\]

Completando cuadrados:

\[
(x-1)^2+(y+1)^2=5.
\]

Por tanto:

\[
\boxed{
\text{Circunferencia de centro }(1,-1)
\text{ y radio }\sqrt5.
}
\]

\section{Ejercicio 9}
\subsection{Solución}

La cónica es:

\[
x^2-2\alpha xy+2\alpha y^2-2x+4\alpha y=0.
\]

Se tiene:

\[
|A|=-2\alpha,
\]

\[
|A_0|
=
-\alpha(\alpha-2).
\]

Si:

\[
\alpha<0
\quad\text{o}\quad
\alpha>2,
\]

entonces:

\[
|A_0|<0
\]

y la cónica es una hipérbola.

Si:

\[
0<\alpha<2,
\]

entonces:

\[
|A_0|>0
\]

y:

\[
|A|\tr(A_0)<0.
\]

Por tanto, es una elipse real.

Para:

\[
\alpha=2,
\]

se tiene:

\[
|A_0|=0,
\qquad
|A|\neq0,
\]

por lo que es una parábola.

Para:

\[
\alpha=0,
\]

la ecuación queda:

\[
x^2-2x=0.
\]

Así:

\[
x(x-2)=0.
\]

Por tanto:

\[
\boxed{
\begin{array}{c|c}
\alpha & \text{cónica}\\
\hline
\alpha<0 & \text{hipérbola}\\
\alpha=0 & \text{dos rectas paralelas }x=0,\ x=2\\
0<\alpha<2 & \text{elipse real}\\
\alpha=2 & \text{parábola}\\
\alpha>2 & \text{hipérbola}
\end{array}
}
\]

\section{Ejercicio 10}
\subsection{Solución}

La cónica es:

\[
F(x,y)
=
9x^2+4y^2-18x-16y-11=0.
\]

En:

\[
P=(-1,2),
\]

se tiene:

\[
F_x=-36,
\qquad
F_y=0.
\]

Por tanto, la tangente es:

\[
\boxed{x=-1.}
\]

Para las tangentes desde:

\[
Q=(5,5),
\]

tomamos:

\[
y-5=m(x-5).
\]

Sustituyendo en la cónica y anulando el discriminante se obtiene:

\[
m(m-2)=0.
\]

Así:

\[
\boxed{
y=5
}
\]

y:

\[
\boxed{
y=2x-5.
}
\]

\section{Ejercicio 11}
\subsection{Solución}

La cónica es:

\[
8x^2+2y^2-4y+2=0.
\]

Su matriz homogénea es:

\[
A=
\begin{pmatrix}
8&0&0\\
0&2&-2\\
0&-2&2
\end{pmatrix}.
\]

Para:

\[
P=(0,1),
\]

la ecuación polar es:

\[
\begin{pmatrix}
0&1&1
\end{pmatrix}
A
\begin{pmatrix}
x\\y\\1
\end{pmatrix}
=0.
\]

Pero:

\[
\begin{pmatrix}
0&1&1
\end{pmatrix}
A
=
\begin{pmatrix}
0&0&0
\end{pmatrix}.
\]

Por tanto:

\[
\boxed{
0=0.
}
\]

La cónica es degenerada:

\[
8x^2+2(y-1)^2=0,
\]

por lo que representa únicamente el punto \((0,1)\). En consecuencia, en este caso no se obtiene una recta polar no trivial.

\section{Ejercicio 12}
\subsection{Solución}

La cónica es:

\[
x^2-4xy+y^2-6x+2y=0.
\]

Las rectas que pasan por:

\[
P=(-1,-3)
\]

son:

\[
y+3=m(x+1).
\]

Sustituyendo en la cónica y anulando el discriminante:

\[
(m-1)(m-3)=0.
\]

Por tanto:

\[
\boxed{
y=x-2
}
\]

y:

\[
\boxed{
y=3x.
}
\]

\section{Ejercicio 13}
\subsection{Solución}

La hipérbola es:

\[
6xy-2x-1=0.
\]

Su centro satisface:

\[
A_0C=-L^t.
\]

Se obtiene:

\[
\boxed{
C=
\left(
0,\frac13
\right).
}
\]

La ecuación puede escribirse como:

\[
6x\left(
y-\frac13
\right)-1=0.
\]

Por tanto, sus asíntotas son:

\[
\boxed{
x=0
}
\]

y:

\[
\boxed{
y=\frac13.
}
\]

\section{Ejercicio 14}
\subsection{Solución}

La cónica es:

\[
3x^2-5y^2-2xy-4x+2
+
\alpha(x^2-2y^2-xy)
=
0.
\]

Para que la hipérbola sea equilátera debe cumplirse:

\[
\tr(A_0)=0.
\]

Así:

\[
3+\alpha-5-2\alpha=0.
\]

Por tanto:

\[
\boxed{
\alpha=-2.
}
\]

La ecuación queda:

\[
x^2-y^2-4x+2=0.
\]

Completando cuadrados:

\[
(x-2)^2-y^2=2.
\]

Las asíntotas son:

\[
\boxed{
y=x-2
}
\]

y:

\[
\boxed{
y=-x+2.
}
\]

\section{Ejercicio 15}
\subsection{Solución}

La hipérbola es:

\[
\frac{x^2}{4}-y^2=1.
\]

Sea:

\[
P=(x_0,y_0)
\]

un punto de la hipérbola.

Un vector normal viene dado por el gradiente:

\[
\nabla F(P)
=
\left(
\frac{x_0}{2},
-2y_0
\right).
\]

Por tanto, la normal en \(P\) es:

\[
\boxed{
(x,y)
=
(x_0,y_0)
+
\lambda
\left(
\frac{x_0}{2},
-2y_0
\right).
}
\]

Equivalentemente:

\[
\boxed{
y-y_0
=
-\frac{4y_0}{x_0}(x-x_0).
}
\]

\section{Ejercicio 16}
\subsection{Solución}

La cónica es:

\[
11x^2+14y^2-4xy+40x+20y+45=0.
\]

La matriz cuadrática es:

\[
A_0=
\begin{pmatrix}
11&-2\\
-2&14
\end{pmatrix}.
\]

Sus autovalores son:

\[
\lambda_1=10,
\qquad
\lambda_2=15,
\]

con autovectores ortonormales:

\[
u_1=
\frac1{\sqrt5}(2,1),
\qquad
u_2=
\frac1{\sqrt5}(-1,2).
\]

El centro satisface:

\[
A_0C=-L^t.
\]

Así:

\[
\boxed{
C=(-2,-1).
}
\]

Tomando coordenadas \((X,Y)\) respecto de los ejes \(u_1,u_2\) con origen \(C\), se obtiene:

\[
10X^2+15Y^2-5=0.
\]

Por tanto:

\[
\boxed{
\frac{X^2}{1/2}
+
\frac{Y^2}{1/3}
=
1.
}
\]

Así, la cónica es una elipse real.

Se tiene:

\[
a=\frac1{\sqrt2},
\qquad
b=\frac1{\sqrt3},
\]

\[
c=
\sqrt{a^2-b^2}
=
\frac1{\sqrt6}.
\]

La excentricidad es:

\[
\boxed{
e=\frac1{\sqrt3}.
}
\]

Los ejes son:

\[
\boxed{
(x,y)
=
(-2,-1)
+
t(2,1)
}
\]

y:

\[
\boxed{
(x,y)
=
(-2,-1)
+
t(-1,2).
}
\]

Los cuatro vértices son:

\[
\boxed{
C\pm
\frac1{\sqrt{10}}(2,1)
}
\]

y:

\[
\boxed{
C\pm
\frac1{\sqrt{15}}(-1,2).
}
\]

Los focos son:

\[
\boxed{
F_{1,2}
=
(-2,-1)
\pm
\frac1{\sqrt{30}}(2,1).
}
\]

\section{Ejercicio 17}
\subsection{Solución}

La cónica es:

\[
x^2-4xy+4y^2-2x+8y=0.
\]

La matriz cuadrática es:

\[
A_0=
\begin{pmatrix}
1&-2\\
-2&4
\end{pmatrix}.
\]

Sus autovalores son:

\[
0,
\qquad
5.
\]

Tomamos:

\[
u=
\frac1{\sqrt5}(1,-2),
\qquad
v=
\frac1{\sqrt5}(2,1).
\]

Con:

\[
\begin{pmatrix}
x\\y
\end{pmatrix}
=
u\,\xi+v\,\eta,
\]

se obtiene:

\[
5\xi^2
-\frac{18\sqrt5}{5}\xi
+
\frac{4\sqrt5}{5}\eta
=
0.
\]

Definimos:

\[
\xi'
=
\xi-\frac{9\sqrt5}{25},
\]

\[
\eta'
=
\eta-\frac{81}{20\sqrt5}.
\]

Entonces:

\[
\boxed{
(\xi')^2
=
-\frac{4\sqrt5}{25}\eta'.
}
\]

Por tanto, es una parábola y no tiene centro.

El vértice es:

\[
\boxed{
V=
\left(
\frac{99}{50},
\frac9{100}
\right).
}
\]

Su eje tiene dirección:

\[
(2,1)
\]

y ecuación:

\[
\boxed{
x-2y=\frac95.
}
\]

Como:

\[
4p=-\frac{4\sqrt5}{25},
\]

se tiene:

\[
p=-\frac{\sqrt5}{25}.
\]

La excentricidad es:

\[
\boxed{e=1.}
\]

El foco es:

\[
\boxed{
F=
\left(
\frac{19}{10},
\frac1{20}
\right).
}
\]

\section{Ejercicio 18}
\subsection{Solución}

La cónica es:

\[
x^2-8xy+y^2-4x-4y+1=0.
\]

La matriz cuadrática es:

\[
A_0=
\begin{pmatrix}
1&-4\\
-4&1
\end{pmatrix}.
\]

Sus autovalores son:

\[
-3,
\qquad
5,
\]

con autovectores ortonormales:

\[
u=
\frac1{\sqrt2}(1,1),
\qquad
v=
\frac1{\sqrt2}(1,-1).
\]

El centro es:

\[
\boxed{
C=
\left(
-\frac23,-\frac23
\right).
}
\]

En coordenadas \((X,Y)\) respecto de \(u,v\):

\[
-3X^2+5Y^2+\frac{11}{3}=0.
\]

Por tanto:

\[
\boxed{
\frac{X^2}{11/9}
-
\frac{Y^2}{11/15}
=
1.
}
\]

Es una hipérbola.

Se tiene:

\[
a=\frac{\sqrt{11}}{3},
\qquad
b=\sqrt{\frac{11}{15}},
\]

\[
c=\sqrt{a^2+b^2},
\qquad
e=\frac ca.
\]

Así:

\[
\boxed{
e=\sqrt{\frac85}.
}
\]

Los ejes son:

\[
\boxed{
y=x
}
\]

y:

\[
\boxed{
x+y=-\frac43.
}
\]

Los vértices son:

\[
\boxed{
\left(
-\frac23\pm\frac{\sqrt{22}}6,
-\frac23\pm\frac{\sqrt{22}}6
\right).
}
\]

Los focos son:

\[
\boxed{
\left(
-\frac23\pm\frac{2\sqrt{55}}{15},
-\frac23\pm\frac{2\sqrt{55}}{15}
\right).
}
\]

\section{Ejercicio 19}
\subsection{Solución}

La cónica es:

\[
4x^2+\alpha y^2+2x-2=0.
\]

Completando cuadrados:

\[
4\left(
x+\frac14
\right)^2
+
\alpha y^2
=
\frac94.
\]

Si:

\[
\alpha>0,
\]

se obtiene:

\[
\boxed{
\frac{(x+\frac14)^2}{9/16}
+
\frac{y^2}{9/(4\alpha)}
=
1.
}
\]

Es una elipse real.

Si:

\[
\alpha=0,
\]

se tiene:

\[
4\left(
x+\frac14
\right)^2
=
\frac94.
\]

Por tanto:

\[
\boxed{
x=\frac12
\quad\text{o}\quad
x=-1.
}
\]

Si:

\[
\alpha<0,
\]

se obtiene:

\[
\boxed{
\frac{(x+\frac14)^2}{9/16}
-
\frac{y^2}{-9/(4\alpha)}
=
1.
}
\]

Es una hipérbola.

\section{Ejercicio 20}
\subsection{Solución}

La familia es:

\[
x^2+m^2y^2+2xy+2x+2my+2m=0.
\]

Se tiene:

\[
|A|
=
2m^2(m-1),
\]

\[
|A_0|
=
(m-1)(m+1).
\]

Por tanto:

\[
\boxed{
\begin{array}{c|c}
m & \text{cónica}\\
\hline
m<-1 & \text{elipse real}\\
m=-1 & \text{parábola}\\
-1<m<1,\ m\neq0 & \text{hipérbola}\\
m=0 & \text{dos rectas secantes}\\
m=1 & \text{dos rectas paralelas imaginarias}\\
m>1 & \text{elipse imaginaria}
\end{array}
}
\]

Para:

\[
m=-1,
\]

la ecuación es:

\[
(x+y)^2+2(x-y)-2=0.
\]

Tomamos:

\[
\xi=\frac{x+y}{\sqrt2},
\qquad
\eta=\frac{x-y}{\sqrt2}.
\]

Entonces:

\[
2\xi^2+2\sqrt2\,\eta-2=0.
\]

Definiendo:

\[
\eta'
=
\eta-\frac1{\sqrt2},
\]

se obtiene:

\[
\boxed{
\xi^2=-\sqrt2\,\eta'.
}
\]

El vértice es:

\[
V=
\left(
\frac12,-\frac12
\right).
\]

Como:

\[
4p=-\sqrt2,
\]

se tiene:

\[
p=-\frac1{2\sqrt2}.
\]

El foco es:

\[
\boxed{
F=
\left(
\frac14,-\frac14
\right).
}
\]

\section{Ejercicio 21}
\subsection{Solución}

La cónica es:

\[
3y^2+4xy-4x-6y-1=0.
\]

La matriz cuadrática es:

\[
A_0=
\begin{pmatrix}
0&2\\
2&3
\end{pmatrix}.
\]

Se tiene:

\[
|A_0|=-4<0.
\]

Por tanto:

\[
\boxed{\text{Es una hipérbola}.}
\]

El centro es:

\[
\boxed{
C=(0,1).
}
\]

Los autovalores son:

\[
-1,
\qquad
4,
\]

con autovectores ortonormales:

\[
u=
\frac1{\sqrt5}(-2,1),
\qquad
v=
\frac1{\sqrt5}(1,2).
\]

En coordenadas \((X,Y)\) respecto de \(u,v\):

\[
-X^2+4Y^2-4=0.
\]

Por tanto:

\[
\boxed{
Y^2-\frac{X^2}{4}=1.
}
\]

Así:

\[
a=1,
\qquad
b=2,
\qquad
c=\sqrt5.
\]

Los focos son:

\[
\boxed{
F_1=(1,3),
\qquad
F_2=(-1,-1).
}
\]

Para obtener el punto cuya polar es:

\[
x=1,
\]

tomamos:

\[
P=(x_0,y_0)
\]

y exigimos que:

\[
(x_0,y_0,1)A
\]

sea proporcional a:

\[
(1,0,-1).
\]

Se obtiene:

\[
\boxed{
P=(-3,3).
}
\]

\section{Ejercicio 22}
\subsection{Solución}

\textbf{a)}

\[
4x^2+z^2+4xz-4y+1=0.
\]

La parte cuadrática es:

\[
(2x+z)^2.
\]

Se tiene:

\[
\rg(A)=3,
\qquad
\rg(A_0)=1.
\]

Por tanto:

\[
\boxed{\text{Cilindro parabólico}.}
\]

\textbf{b)}

\[
x^2+y^2+9z^2-6xz+2x-2y+2=0.
\]

La matriz \(A_0\) es semidefinida positiva de rango \(2\), mientras que:

\[
|A|<0.
\]

Por tanto:

\[
\boxed{\text{Paraboloide elíptico}.}
\]

\textbf{c)}

\[
xy+yz+xz-1=0.
\]

La matriz cuadrática tiene autovalores:

\[
1,
\qquad
-\frac12,
\qquad
-\frac12.
\]

Por tanto:

\[
\boxed{\text{Hiperboloide elíptico de dos hojas}.}
\]

\section{Ejercicio 23}
\subsection{Solución}

La cuádrica es:

\[
9y^2+4z^2-36=0.
\]

Dividiendo entre \(36\):

\[
\boxed{
\frac{y^2}{4}
+
\frac{z^2}{9}
=
1.
}
\]

La variable \(x\) es libre.

Por tanto:

\[
\boxed{\text{Cilindro elíptico}.}
\]

\section{Ejercicio 24}
\subsection{Solución}

La cuádrica es:

\[
x^2+y^2-5z^2+4xy-8xz+8yz
+18x-18y+36z-54=0.
\]

La matriz cuadrática es:

\[
A_0=
\begin{pmatrix}
1&2&-4\\
2&1&4\\
-4&4&-5
\end{pmatrix}.
\]

Sus autovalores son:

\[
-9,
\qquad
3,
\qquad
3.
\]

El centro es:

\[
\boxed{
C=(1,-1,2).
}
\]

Tomamos la base ortonormal:

\[
u_1=
\frac1{\sqrt6}(1,-1,2),
\]

\[
u_2=
\frac1{\sqrt2}(1,1,0),
\]

\[
u_3=
\frac1{\sqrt3}(1,-1,-1).
\]

En coordenadas \((X,Y,Z)\) con origen \(C\):

\[
-9X^2+3Y^2+3Z^2=0.
\]

Por tanto:

\[
\boxed{
Y^2+Z^2=3X^2.
}
\]

Así:

\[
\boxed{\text{Cono elíptico real}.}
\]

\section{Ejercicio 25}
\subsection{Solución}

La cuádrica es:

\[
11x^2+11y^2+14z^2
-2xy-8xz-8yz
+6x-18y-60z+104=0.
\]

La matriz cuadrática tiene autovalores:

\[
6,
\qquad
12,
\qquad
18.
\]

El centro es:

\[
\boxed{
C=(1,2,3).
}
\]

Tomamos la base ortonormal:

\[
u_1=
\frac1{\sqrt3}(1,1,1),
\]

\[
u_2=
\frac1{\sqrt2}(-1,1,0),
\]

\[
u_3=
\frac1{\sqrt6}(-1,-1,2).
\]

En coordenadas \((X,Y,Z)\) con origen \(C\):

\[
6X^2+12Y^2+18Z^2-1=0.
\]

Por tanto:

\[
\boxed{
6X^2+12Y^2+18Z^2=1.
}
\]

Así:

\[
\boxed{\text{Elipsoide real}.}
\]

\section{Ejercicio 26}
\subsection{Solución}

La cuádrica es:

\[
x^2+y^2+z^2-2xy+2xz+3x-y+z+1=0.
\]

La matriz cuadrática es:

\[
A_0=
\begin{pmatrix}
1&-1&1\\
-1&1&0\\
1&0&1
\end{pmatrix}.
\]

El centro satisface:

\[
A_0C=-L^t.
\]

Se obtiene:

\[
\boxed{
C=
\left(
\frac12,1,-1
\right).
}
\]

Los autovalores son:

\[
1,
\qquad
1-\sqrt2,
\qquad
1+\sqrt2,
\]

con autovectores:

\[
v_1=(0,1,1),
\]

\[
v_2=(-\sqrt2,-1,1),
\]

\[
v_3=(\sqrt2,-1,1).
\]

Por tanto, los ejes son:

\[
\boxed{
r_1:
X=C+\lambda(0,1,1)
}
\]

\[
\boxed{
r_2:
X=C+\lambda(-\sqrt2,-1,1)
}
\]

\[
\boxed{
r_3:
X=C+\lambda(\sqrt2,-1,1).
}
\]

\section{Ejercicio 27}
\subsection{Solución}

La cuádrica es:

\[
x^2+2y^2+z^2-2xz+4x+4y+5=0.
\]

La matriz cuadrática tiene autovalores:

\[
0,
\qquad
2,
\qquad
2.
\]

Tomamos:

\[
u_0=
\frac1{\sqrt2}(1,0,1),
\]

\[
u_1=
\frac1{\sqrt2}(1,0,-1),
\]

\[
u_2=(0,1,0).
\]

Con coordenadas \((\xi,\eta,\zeta)\) respecto de esta base:

\[
2\eta^2+2\zeta^2
+
2\sqrt2\,\xi
+
2\sqrt2\,\eta
+
4\zeta
+
5
=
0.
\]

Completando cuadrados:

\[
\eta'=
\eta+\frac1{\sqrt2},
\]

\[
\zeta'=\zeta+1,
\]

\[
\xi'=
\xi+\frac1{\sqrt2},
\]

se obtiene:

\[
\boxed{
(\eta')^2+(\zeta')^2+\sqrt2\,\xi'=0.
}
\]

Por tanto, el vértice es:

\[
\boxed{
V=(-1,-1,0).
}
\]

El eje principal tiene dirección:

\[
(1,0,1).
\]

Así:

\[
\boxed{
r:
X=(-1,-1,0)+\lambda(1,0,1).
}
\]

\section{Ejercicio 28}
\subsection{Solución}

La cuádrica es:

\[
x^2+2\alpha xy+2xz+z^2-2x+4y-z+1=0.
\]

La matriz cuadrática es:

\[
A_0=
\begin{pmatrix}
1&\alpha&1\\
\alpha&0&0\\
1&0&1
\end{pmatrix}.
\]

Su determinante es:

\[
|A_0|=-\alpha^2.
\]

Si:

\[
\alpha\neq0,
\]

existe un centro único.

Resolviendo:

\[
A_0C=-L^t,
\]

se obtiene:

\[
\boxed{
C=
\left(
-\frac2\alpha,
\frac1{2\alpha},
\frac{\alpha+4}{2\alpha}
\right).
}
\]

Para:

\[
\alpha=0,
\]

el sistema que determina el centro es incompatible.

Por tanto:

\[
\boxed{
\text{La cuádrica tiene centro si y solo si }\alpha\neq0.
}
\]

\section{Ejercicio 29}
\subsection{Solución}

El hiperboloide es:

\[
\frac{x^2}{4}
+
\frac{y^2}{9}
-
\frac{z^2}{16}
=
1.
\]

Sea:

\[
F(x,y,z)
=
\frac{x^2}{4}
+
\frac{y^2}{9}
-
\frac{z^2}{16}
-
1.
\]

En:

\[
P=(2,0,0),
\]

se tiene:

\[
\nabla F(P)=(1,0,0).
\]

Por tanto, el plano tangente es:

\[
\boxed{
x=2.
}
\]

Intersecando con el hiperboloide:

\[
1+\frac{y^2}{9}-\frac{z^2}{16}=1.
\]

Así:

\[
\frac{y^2}{9}
=
\frac{z^2}{16}.
\]

Por tanto, la intersección está formada por las dos rectas:

\[
\boxed{
\begin{cases}
x=2,\\
y=\frac34z,
\end{cases}
}
\]

y:

\[
\boxed{
\begin{cases}
x=2,\\
y=-\frac34z.
\end{cases}
}
\]

\section{Ejercicio 30}
\subsection{Solución}

La cuádrica es:

\[
Q:
x^2+y^2+2xy-2z=0.
\]

Su matriz homogénea es:

\[
A=
\begin{pmatrix}
1&1&0&0\\
1&1&0&0\\
0&0&0&-1\\
0&0&-1&0
\end{pmatrix}.
\]

Tomamos:

\[
P=(0,1,-1).
\]

Se tiene:

\[
\varphi(P,P)=3,
\]

\[
\varphi(X,X)
=
x^2+2xy+y^2-2z,
\]

\[
\varphi(X,P)
=
x+y-z+1.
\]

El cono tangente viene dado por:

\[
\varphi(P,P)\varphi(X,X)
-
\varphi(X,P)^2
=
0.
\]

Por tanto:

\[
\boxed{
3(x^2+2xy+y^2-2z)
-
(x+y-z+1)^2
=
0.
}
\]

Equivalentemente:

\[
\boxed{
2x^2+4xy+2xz-2x
+
2y^2+2yz-2y
-z^2-4z-1
=
0.
}
\]

El contorno aparente verifica:

\[
\varphi(X,X)=0
\]

y:

\[
\varphi(X,P)=0.
\]

Por tanto:

\[
\boxed{
\begin{cases}
x^2+2xy+y^2-2z=0,\\
x+y-z+1=0.
\end{cases}
}
\]

\end{document}
